\clearpage
\chapter*{ЗАКЛЮЧЕНИЕ}
\addcontentsline{toc}{chapter}{ЗАКЛЮЧЕНИЕ}

В ходе выполнения выпускной квалификационной работы была разработана
вычислительная модель расчёта термодинамических свойств воды и водяного пара,
основанная на стандарте IAPWS-IF97 и дополненная численной схемой на основе
метода секущих для маршрута расчёта по параметрам плотности и температуры.
Поставленная цель создания единого программного комплекса для инженерных
расчётов, настольного использования и сервисной интеграции достигнута.

В теоретической части работы сформулирована математическая основа расчёта,
описаны регионы IF97, линия насыщения и принципы получения термодинамических
свойств через производные безразмерных потенциалов.
В практической части реализовано вычислительное ядро на языке Rust с единым API,
типобезопасными единицами измерения и механизмами валидации входных данных.
На базе ядра разработаны настольный интерфейс на FLTK, кроссплатформенный
desktop-вариант на Yew и Tauri, а также gRPC-сервис для пакетных вычислений.

Дополнительно подготовлены средства контейнеризации и развёртывания в среде
Kubernetes, а также сценарии CI/CD для автоматизированной проверки и публикации
артефактов.
Результаты тестирования подтверждают пригодность полученного решения для
практического применения в инженерных и сервисных сценариях.

Дальнейшее развитие работы может быть связано с расширением набора поддерживаемых
диаграмм, углублением средств визуализации, добавлением новых маршрутов расчёта
и подготовкой расширенного набора верификационных материалов для приложений к
диплому.
